\documentstyle[% 


righttag% 


,11pt% 


,amssymb% 


,russian%               remove in Eng. 


,izvru%                 Set izven% in Eng. 


% ,izvru-ks             for brief communications 


]{amsart} 


 


%       Math definitions 


 


\newcommand{\const}{\operatorname{const}} 


\newcommand{\sign}{\operatorname{sign}} 


\newcommand{\ctg}{\operatorname{ctg}} 


\renewcommand{\baselinestretch}{0.96}


\newcommand{\tts}[1]{} 


 


%\hoffset=-15mm%                  For Eng. 


  


\setcounter{page}{3} 


\god{1998} 


\nomer{10~(437)} %                        Remove in Eng. 


%\nomer{Vol.\,42, No.\,10}%               Set in Eng. 


%\str{\pageref{begin}--\pageref{end}}%  Set in Eng. 


\udk{517.546} 


\author{\it И.А.\,Александров, Г.Д.\,Садритдинова} 


% NB:   ^^ remove in Eng. 


\title{ Отображения с симметрией вращения} 


 


\date{} 


% \pagestyle{myheadings}%         Set in Eng. 


 


\pagestyle{plain} %              Remove in Eng. 


\begin{document} 


 


% \thispagestyle{plain}%          Set in Eng. 


 


\maketitle 


\label{begin} 


 


{\it Введение.} Пусть $S_p(M)$ --- множество всех голоморфных в круге 


$E=\{z:|z|<1\}$ 


функций $f(z)=z+\dotsb$, однолистно отображающих $E$ на области из круга 


$\{w:|w|<M\}$ ($1<M\le\infty$), имеющие $p$-кратную ($p\in\Bbb N$)


симметрию вращения относительно точки $w=0$. 


 


Плотный в $S_p(M)$ подкласс функций (в топологии равномерной сходимости 


внутри $E$) образуют функции $w=M\zeta(z,\ln M)$, где $\zeta(z,\tau)$, 


$0\le\tau\le\ln M$, --- решение уравнения Л\"евнера (\cite{1}, с.\,27) 


\begin{equation} 


\frac{d\zeta}{d\tau}= 


-\zeta\frac{\mu^p(\tau)+\zeta^p}{\mu^p(\tau)-\zeta^p},\quad 


\zeta(z,0)=z\in E, 


\end{equation} 


в котором $\mu(\tau)$ --- непрерывная функция с модулем, равным единице. 


Аналогичное предложение справедливо для класса 


$S_p(\infty)\equiv S_p$. 


 


Примеры интегрирования уравнения (1) в квадратурах единичны. Кроме 


простейшего случая $\mu(\tau)=e^{i\alpha}$, $\alpha=\const$, отметим 


примеры, рассмотренные в \cite{2}--\cite{4}. 


 


В данной работе указывается управляющая функция $\mu(\tau)$, для 


которой уравнение (1) приводит к формуле для $\zeta(z,\tau)$


более общей, чем в \cite{2}--\cite{4}. 


\medskip 


 


{\it Выбор $\mu(\tau)$.} Пусть в (1) 


\begin{equation} 


\mu(\tau)=e^{-(\alpha+\delta\tau)i} 


(\lambda(\tau))^{1+\frac{1}{2s}}, 


\end{equation} 


где $\alpha$, $\delta$ --- вещественные постоянные, 


$s$ --- положительное 


число и $\lambda(\tau)$, 


$\lambda(0)=1$, --- некоторая непрерывно дифференцируемая функция на 


$(0,\ln M)$ с модулем, равным единице. Выполнив в (1) замену 


\begin{equation} 


\omega=e^{(\alpha+\delta\tau)i}(\lambda(\tau))^{-\frac{1}{2s}}\zeta, 


\quad \omega(0)=e^{i\alpha}z, 


\end{equation} 


получим для $\omega$ уравнение 


$$ 


\frac{\omega'}{\omega}-i\delta+\frac{1}{2s}\frac{\lambda'}{\lambda}= 


-\frac{\lambda^p+\omega^p}{\lambda^p-\omega^p} 


$$ 


или, что то же самое, 


\begin{equation} 


\frac{\omega'}{\omega}=\bigg(1+i\delta-\frac{1}{2s} 


\frac{\lambda'}{\lambda}\bigg) 


\frac{\eta^p-\omega^p}{\lambda^p-\omega^p}, 


\end{equation} 


где 


$$ 


\eta^p=\eta^p(\tau)=\lambda^p 


\frac{\lambda'+2s(1-i\delta)\lambda}{\lambda'-2s(1+i\delta)\lambda}. 


$$ 


 


Положим 


$$ 


\lambda(\tau)=e^{2i\psi(\tau)},\quad \psi(0)=0. 


$$ 


Функция 


$$ 


\eta^p(\tau)=e^{2ip\psi(\tau)} 


\frac{s+i(\psi'(\tau)-\delta)}{-s+i(\psi'(\tau)-\delta)} 


$$ 


имеет модуль, равный единице. В дальнейшем рассматривается только тот 


случай, когда 


$$ 


\eta^p(\tau)=e^{2ip\varphi}=\gamma^p 


$$ 


--- постоянная. При этом $\psi(\tau)$, как легко показать, является 


решением уравнения 


\begin{equation} 


\psi'=s(\delta+\ctg(\psi-\varphi)p),\quad 


\psi(0)=0. 


\end{equation} 


 


Интегрирование уравнения, эквивалентного (5), 


$$ 


\lambda^p\frac{\lambda'+2\ov m\lambda}{\lambda'-2m\lambda}= 


\gamma^p,\quad \lambda(0)=1, 


$$ 


где $m=s(1+i\delta)$, т.\,е.~уравнения 


$$ 


\frac{d\lambda}{d\tau}= 


-\frac{2\ov m\lambda(\lambda^p+\varepsilon\gamma^p)} 


{\lambda^p-\gamma^p},\quad \lambda(0)=1, 


$$ 


в котором $\varepsilon=m/\ov m$, приводит к формуле 


\begin{equation} 


\lambda^p+\varepsilon\gamma^p=(1+\varepsilon\gamma^p) 


\lambda^{\frac{p}{1+\varepsilon}} 


e^{-\frac{2mp\tau}{1+\varepsilon}} 


\end{equation} 


для нахождения $\lambda(\tau)$, а следовательно, и функции $\mu(\tau)$ ---


по формуле (2). 


 


Из (6) следует, что при $M=\infty$ 


\begin{equation} 


\lim_{\tau\to\infty}\lambda^p(\tau)=-\varepsilon\gamma^p. 


\end{equation} 


 


{\it Функция $\psi(\tau)$.} В (5) произведем замены 


$$ 


(\psi-\varphi)p=y,\quad ps\tau=t 


$$ 


и придем к уравнению 


\begin{equation} 


\frac{dy}{dt}=\delta+\ctg y. 


\end{equation} 


Общий интеграл уравнения (8) дается формулой 


$$ 


\delta y-\ln(\cos y+\delta\sin y)= 


(1+\delta^2)t+\const, 


$$ 


из которой с учетом условия $y(0)=-p\varphi/2$ получаем уравнение 


$$ 


\delta(y+p\varphi)-\ln 


\frac{\cos y+\delta\sin y}{\cos p\varphi-\delta\sin p\varphi} 


=(1+\delta^2)t, 


$$ 


в явной форме задающее функцию $t=t(y)$, $t(-p\varphi)=0$. 


 


Если $\delta=0$, то 


$$ 


\cos y=\cos p\varphi\cdot e^{-t} 


$$ 


и, значит, соответствующая управляющая функция $\mu(\tau)$ имеет вид 


$$ 


\mu(\tau)=e^{-i\alpha}(\lambda(\tau))^{1+\frac{1}{2s}}, 


$$ 


где 


\begin{equation} 


\lambda^p(\tau)=e^{2ip\varphi}\big[\cos p\varphi\cdot e^{-t} 


-i\sqrt{1-\cos^2p\varphi\cdot e^{-2t}}\,\big]^2. 


\end{equation} 


 


Если $\delta\ne 0$, то, полагая $\psi-s\delta\tau=u$, 


$\varphi-s\delta\tau=v$, видим из (5), что 


$$ 


\sign\frac{u'}{s}=\sign(\ctg(u-v)p) 


$$ 


и, следовательно, функция $u$ монотонно изменяется, если ее значения 


принадлежат любому из промежутков 


$$ 


k\pi/2<(u-v)p<(k+1)\pi/2\quad (k\in\Bbb Z). 


$$ 


 


{\it Интегрирование уравнения} (4). Учитывая равенство 


$$ 


1+i\delta-\frac{\lambda'}{2s\lambda}= 


\frac{2\lambda^p}{\lambda^p-\gamma^p}, 


$$ 


представим уравнение (4) в виде 


$$ 


\frac{d\omega}{d\tau}= 


\frac{2\lambda^p\omega(\gamma^p-\omega^p)} 


{(\lambda^p-\gamma^p)(\lambda^p-\omega^p)},\quad 


\omega(0)=z_1=e^{i\alpha}z. 


$$ 


Перейдем от переменной $\tau$ к $\lambda$. Получим 


$$ 


\frac{d\omega}{d\lambda}=- 


\frac{\lambda^{p-1}\omega(\gamma^p-\omega^p)} 


{\ov m(\lambda^p-\omega^p)(\lambda^p+\varepsilon\gamma^p)},\quad 


\omega\big|_{\lambda=1}=z_1, 


$$ 


где $m=s(1+i\delta)$, $\varepsilon=m/\ov m$. Выполнив в этом уравнении 


замену переменной $\lambda$ на $v$ по формуле 


$$ 


v=\frac{1}{\lambda^p+\varepsilon\gamma^p}, 


$$ 


будем иметь 


$$ 


\frac{d\omega}{dv}= 


\frac{\omega(\gamma^p-\omega^p)} 


{\ov m pv\big(\frac{1}{v}-\varepsilon\gamma^p-\omega^p\big)},\quad 


\omega\big|_{v=\frac{1}{1+\varepsilon\gamma^p}}=z_1. 


$$ 


Нетрудно заметить, что производная $dv/d\omega$ линейна относительно 


$v$. Проинтегрируем уравнение 


$$ 


\frac{dv}{d\omega}=- 


\frac{\ov m p(\varepsilon\gamma^p+\omega^p)} 


{\omega(\gamma^p-\omega^p)}v+ 


\frac{\ov m p}{\omega(\gamma^p-\omega^p)},\quad 


v\big|_{\omega=z_1}=\frac{1}{1+\varepsilon\gamma^p}, 


$$ 


варьируя произвольную постоянную $C$ в общем решении 


$$ 


v(\omega)=\frac{C(\gamma^p-\omega^p)^{2s}}{\omega^{mp}} 


$$ 


соответствующего однородного уравнения. С учетом начального условия 


получим 


$$ 


\frac{\omega^{mp}}{(\gamma^p-\omega^p)^{2s}}v=\ov m p 


\int^\omega_{z_1}\frac{u^{mp-1}du}{(\gamma^p-u^p)^{2s+1}}- 


\frac{z^{mp}_1}{(1+\varepsilon\gamma^p)(\gamma^p-z^p_1)^{2s}}. 


$$ 


Так как 


$$ 


\frac{z^{mp}_1}{(\gamma^p-z^p_1)^{2s}}=\ov m p\int^{z_1}_0 


\frac{u^{mp-1}(\varepsilon\gamma^p+u^p)du} 


{(\gamma^p-u^p)^{2s+1}}, 


$$ 


то окончательно получим 


\begin{equation} 


\frac{\omega^{mp}v}{\ov mp(\gamma^p-\omega^p)^{2s}}= 


\int^\omega_{z_1}


\frac{u^{mp-1}du}{(\gamma^p-u^p)^{2s+1}}+ 


\frac{1}{1+\varepsilon\gamma^p}\int^{z_1}_0 


\frac{u^{mp-1}(\varepsilon\gamma^p+u^p)du}{(\gamma^p-u^p)^{2s+1}}. 


\end{equation}  


Итак, имеет место 


 


\begin{thm} 


Функция 


$$ 


\zeta(z,\tau)=-e^{-(\alpha+\delta\tau)i}\lambda^{\frac{1}{2s}} 


(\tau)\omega(v(\tau)),\quad 


\zeta(0,\tau)=0,\quad \zeta'_z(0,\tau)>0, 


$$ 


где $\alpha$, $\delta$ --- вещественные постоянные, 


$p\in\Bbb N$, $s\in\Bbb R^+$, $\lambda(\tau)=e^{2i\psi(\tau)}$, причем 


$\psi(\tau)$ определяется уравнением 


$$ 


\delta\psi-\ln 


\frac{\cos(\psi-\varphi)p+\delta\sin(\psi-\varphi)p} 


{\cos p\varphi-\delta\sin p\varphi}= 


(1+\delta^2)ps\tau, 


$$ 


и $\omega(v)$ неявно задается уравнением $(10)$, в котором 


$z_1=e^{i\alpha}z$, $m=s(1+i\delta)$, $\varepsilon=m/\ov m$,


$\gamma=e^{2i\varphi}$ --- 


постоянная, конформно и однолистно отображает круг $E$ на область, 


лежащую в единичном круге и имеющую $p$-кратную симметрию вращения 


относительно нуля. 


\end{thm} 


 


{\it Предельный случай.} При $\tau=\ln M$ функция $M{}\zeta(z,\tau)$ 


входит в класс $S_p(M)$. Устремим $M$ к $\infty$. Предельная функция, 


как известно, голоморфна и однолистна. Найдем ее. Перейдем к пределу 


при $\tau\to\infty$ сначала в правой части 


формулы (10). Имеем 


$$ 


\lim_{\tau\to\infty}\bigg[\,\int^{\omega}_{z_1} 


\frac{u^{mp-1}du}{(\gamma^p-u^p)^{2s+1}}+ 


\frac{1}{1+\varepsilon\gamma^p}\int^{z_1}_0 


\frac{u^{mp-1}(\varepsilon\gamma^p+u^p)du} 


{(\gamma^p-u^p)^{2s+1}}\bigg]=-\frac{1}{1+\varepsilon\gamma^p} 


\int^{z_1}_0 


\frac{(1-u^p)u^{mp-1}du}{(\gamma^p-u^p)^{2s+1}}, 


$$ 


поскольку каждое решение $\zeta(z,\tau)$, $z\in E$, уравнения (1) 


стремится к 


нулю при $\tau\to\infty$ равномерно относительно $z$ внутри $E$ и, 


следовательно, $\lim\limits_{\tau\to\infty}\omega(\tau)=0$. 


 


Выполним теперь предельный переход при $\tau\to\infty$ в левой части 


формулы (10). Учтем, что 


$\lim\limits_{\tau\to\infty}e^\tau\zeta(z,\tau)=f(z)\in S_p$. 


Имеем, используя (6), (7),


\begin{multline*} 


\lim_{\tau\to\infty} 


\frac{\omega^{mp}}{(\lambda^p+\varepsilon\gamma^p) 


(\gamma^p-\omega^p)^{2s}}=\frac{e^{i\alpha mp}}{\gamma^{2ps}} 


\lim_{\tau\to\infty} 


\frac{e^{i\delta\tau mp}\lambda^{-\frac{mp}{2s}}\zeta^{mp}} 


{\lambda^p+\varepsilon\gamma^p} = \\ 


% 


=\frac{e^{i\alpha mp}}{\gamma^{2ps}(1+\varepsilon\gamma^p)} 


\lim_{\tau\to\infty} 


\big[\lambda^{-\frac{mp}{2s}-\frac{p}{1+\varepsilon}}\cdot 


e^{i\delta\tau mp+\frac{2mp\tau}{1+\varepsilon}-\tau mp}\big] 


f^{mp}(z)= \\ 


% 


=\frac{e^{i\alpha mp}}{\gamma^{2ps}(1+\varepsilon\gamma^p)} 


\lim_{\tau\to\infty}\lambda^{-p}\cdot f^{mp}(z)= 


-\frac{\ov m e^{i\alpha mp}f^{mp}(z)} 


{m(1+\varepsilon\gamma^p)\gamma^{(2s+1)p}}. 


\end{multline*} 


 


Приравнивая полученные пределы, найдем 


\begin{equation} 


f(z)=e^{i\alpha}\bigg[ mp\int^{z_1}_0 


\frac{(1-u^p)u^{mp-1}du}{(1-\ov\gamma^{\,p}u^p)^{2s+1}} 


\bigg]^{\frac{1}{mp}}. 


\end{equation} 


 


\begin{thm} 


Пусть $\alpha$, $\delta$ --- вещественные постоянные, 


$p\in\Bbb N$, $s\in\Bbb R^+$, $m=s(1+i\delta)$, $z_1=e^{i\alpha}z$. 


Тогда 


однозначная ветвь функции $f(z)=z+\dotsb$, даваемая формулой $(11)$, 


однолистна в 


единичном круге и отображает его на область с $p$-кратной симметрией 


вращения относительно начала. 


\end{thm} 


 


{\it Случай} $m=1$. Взяв интегралы в (10) при $s=1$, $\delta=0$ и 


$\varepsilon=1$, получаем уравнение 


$$ 


\frac{\omega^p}{(\lambda^p+\gamma^p)(\gamma^p-\omega^p)^2}= 


\frac{1}{2(\gamma^p-\omega^p)^2}+\bigg( 


\frac{z^p_1}{1+\gamma^p}-\frac{1}{2}\bigg) 


\frac{1}{(\gamma^p-z^p_1)^2}, 


$$ 


сводящееся к квадратному алгебраическому уравнению относительно 


$\omega^p$. 


Остановимся здесь на случае, когда $\alpha=\varphi=0$. Простые 


вычисления и формула (3) приводят к соотношению 


$$ 


\zeta^p(z,\tau)=\lambda^{p/2}\omega^p= 


\frac{\lambda^{p/2}}{1+\lambda^p}\big[\lambda^p+z^p-\lambda^p 


\sqrt{(1-z^p)(1-\lambda^{-2p}z^p)}\,\big]. 


$$ 


Функция $\zeta(z,\tau)=e^{-\tau}z+\dotsb$ с $\lambda(\tau)$, 


записываемой по формуле 


(9), отображает единичный круг на круг с $p$ исключенными луночками. 


 


\begin{thebibliography}{9} 


 


\bibitem{1} 


Александров И.А. {\em Параметрические продолжения в теории однолистных 


функций}. -- М.: Наука, 1976. -- 343 с. 


\bibitem{2} 


Куфарев П.П. {\em Одно замечание об интегралах уравнения Л\"евнера} // 


ДАН СССР. -- 1947. -- Т.\,57. -- \No\,7. -- С.\,655--656. 


\bibitem{3} 


Александров И.А. {\em Об одном случае интегрирования уравнения 


Л\"евнера} // Сиб. матем. журн. -- 1981. -- Т.\,22. -- \No\,2. -- 


С.\,207--209. 


\bibitem{4} 


Helling K. {\em Beitr\"age zur Theorie der L\"ownerschen 


Differentialgleichung} // Wiss. Z. P\"ad. Hochsch. 


``Liselotte-Herrmann'' G\"ustrow. Math. -- naturwiss. Fak. -- 1976. -- 


\No\,2. -- S.\,319--328. 


 


\end{thebibliography} 


 


\vskip 0.5cm 


 


\post{\it Томский государственный}{\it Поступила} 


\post{\it университет}{10.04.1996} 


 


\label{end} 


\end{document} 


